\documentclass[11pt]{article}
\usepackage[margin=1in]{geometry}
\usepackage{amsmath,amssymb,mathtools}
\usepackage{microtype}
\usepackage{hyperref}
\usepackage{enumitem}
\usepackage{booktabs}
\setlength{\parindent}{0pt}
\setlength{\parskip}{0.55em}
\title{LRSC Odd-Ring Spectral Rank Theorem\\\large Version 1.2 - Analytic Family Result}
\author{Prince Upadhyay\\Independent Research}
\date{22 September 2026}
\begin{document}
\maketitle

\begin{abstract}
For a single-step near-Nyquist local redistribution map on an odd cyclic lattice, this note proves that the rank of the conjugate-balanced post-stencil channel family equals the number of active reflection orbits, provided at least one reflection orbit is inactive. The proof closes the previously isolated non-vanishing condition for the Fourier coefficients of the local shedding field for every odd $M\ge 9$, using residue-class trigonometric formulas. A Chebyshev/Vandermonde argument gives the pre-stencil rank, and injectivity of the stencil on the masked channel span preserves it. The result is structural and does not imply a universal fixed-unit compression number $K_\tau$.
\end{abstract}

\section{Setup and theorem}
Let $M\ge 9$ be odd and
\[
k=\frac{M-1}{2},\qquad \alpha=\frac{\pi}{M}.
\]
On $\mathbb Z_M$, define the near-Nyquist carrier
\[
X_t=1+A\cos\!\left(\frac{2\pi kt}{M}\right),\qquad A>0,
\]
and the local shedding field
\[
j_t=X_t-\min\{X_{t-1},X_t,X_{t+1}\}.
\]
Let $m_t\in\{0,1\}$ be any reflection-symmetric mask, $m_t=m_{-t}$, with at least one inactive reflection orbit. Let $p$ denote the number of active reflection orbits. Form the conjugate-balanced pre-stencil channel family generated by the Fourier coefficients of $j$, and let $S$ be the circulant stencil whose Fourier multiplier is
\[
s_q=\cos\!\left(\frac{2\pi q}{M}\right)-1.
\]

\textbf{Theorem.} Under these assumptions, the post-stencil channel matrix $V=SW$ satisfies
\[
\boxed{\operatorname{rank}(V)=p.}
\]

\section{Exact shedding field}
Using centered representatives $t=-k,\dots,k$,
\[
\cos\!\left(\frac{2\pi kt}{M}\right)=\cos((\pi-\alpha)t)=(-1)^t\cos(\alpha t),
\]
so
\[
X_t=1+A(-1)^t\cos(\alpha t).
\]
For odd $|t|$, $X_t$ is the local minimum and $j_t=0$. At the origin,
\[
j_0=A(1+\cos\alpha).
\]
For even $|t|>0$, monotonicity of cosine on the centered half-ring gives
\[
\boxed{j_t=A\left[\cos(|t|\alpha)+\cos((|t|-1)\alpha)\right].}
\]
Thus $j$ is real and reflection-even.

Let
\[
L=\left\lfloor\frac{M-1}{4}\right\rfloor.
\]
Then
\[
\widehat j_r=\frac{1}{M}\left[j_0+2\sum_{u=1}^{L}j_{2u}\cos\!\left(\frac{4\pi r u}{M}\right)\right].
\]
Finite trigonometric summation reduces this expression as follows.

\section{Non-vanishing of the shedding spectrum}
\subsection{$M=4L+3$}
For $0\le r<k$,
\[
\widehat j_r=\frac{2A\cos(\alpha/2)}{M}\,B_r,
\]
where
\[
B_r=
-\frac{\sin(2\alpha)\left[(-1)^r\cos(r\alpha)+\sin(\alpha/2)\right]}
{2\sin((2r-1)\alpha)\sin((2r+1)\alpha)}.
\]
The denominator is nonzero for integer $0\le r<k$. For even $r$, the bracket is positive. For odd $r$,
\[
r\alpha\le (k-1)\alpha=\frac{\pi}{2}-\frac{3\alpha}{2},
\]
hence
\[
\cos(r\alpha)\ge \sin(3\alpha/2)>\sin(\alpha/2).
\]
Thus the bracket is strictly negative. Therefore
\[
\boxed{\widehat j_r\ne 0\qquad(0\le r<k)}
\]
for every $M\equiv3\pmod4$.

\subsection{$M=4L+1$}
Here $k=2L$ and
\[
\widehat j_r=\frac{2A\cos(\alpha/2)}{M}\,B_r,
\]
with
\[
B_r=
-\frac{\sin\alpha\left[(-1)^r\cos(r\alpha)\cos(2r\alpha)+\cos\alpha\sin(\alpha/2)\right]}
{\sin((2r-1)\alpha)\sin((2r+1)\alpha)}.
\]
Again the denominator is nonzero for $0\le r<k$. A zero would require
\[
\left|\cos(r\alpha)\cos(2r\alpha)\right|=c,
\qquad c:=\cos\alpha\sin(\alpha/2).
\]

For $0\le r\le L-1$, $f(x)=\cos x\cos2x$ is strictly decreasing on $[0,\pi/4)$. Hence
\[
f(r\alpha)\ge f((L-1)\alpha)
=\cos\!\left(\frac\pi4-\frac{5\alpha}{4}\right)\sin\!\left(\frac{5\alpha}{2}\right).
\]
Since $M\ge9$, $\alpha\le\pi/9$. Therefore
\[
\cos\!\left(\frac\pi4-\frac{5\alpha}{4}\right)>\frac1{\sqrt2},
\qquad
\sin\!\left(\frac{5\alpha}{2}\right)\ge\frac{5\alpha}{\pi},
\]
while $c<\alpha/2$. Because $5/(\pi\sqrt2)>1/2$, one has $f(r\alpha)>c$.

At $r=L$,
\[
\left|\cos(L\alpha)\cos(2L\alpha)\right|
=\cos(L\alpha)\sin(\alpha/2).
\]
Since $L\alpha=\pi/4-\alpha/4>\alpha$, $\cos(L\alpha)<\cos\alpha$, so this quantity is strictly less than $c$.

For $L+1\le r\le2L-1$, put $s=2L-r$ and $t=(s+1/2)\alpha$. Then
\[
\left|\cos(r\alpha)\cos(2r\alpha)\right|=h(t):=\sin t\cos2t,
\]
where
\[
\frac{3\alpha}{2}\le t\le \frac\pi4-\frac{3\alpha}{4}.
\]
Moreover
\[
h'(t)=\cos t\left(1-6\sin^2 t\right),
\]
so $h$ has at most one critical point in this interval, and that point is a maximum. The minimum is therefore attained at an endpoint.

At $t=3\alpha/2$,
\[
\frac{h(3\alpha/2)}{c}
=\frac{\sin(3\alpha/2)}{\sin(\alpha/2)}\frac{\cos3\alpha}{\cos\alpha}
=\frac{(1+2\cos\alpha)\cos3\alpha}{\cos\alpha}>1,
\]
because $\alpha\le\pi/9$ implies $\cos3\alpha\ge1/2$. At the other endpoint,
\[
\frac{h(\pi/4-3\alpha/4)}{c}
=\frac{\sin(\pi/4-3\alpha/4)}{1}\,
\frac{1+2\cos\alpha}{\cos\alpha}>1,
\]
because $\sin(\pi/4-3\alpha/4)\ge1/2$.

Hence equality with $c$ is impossible for every integer $0\le r<k$, and
\[
\boxed{\widehat j_r\ne0\qquad(0\le r<k)}
\]
for every $M\equiv1\pmod4$ as well.

Combining the two residue classes proves the non-vanishing lemma for every odd $M\ge9$.

\section{Pre-stencil rank}
Since $j$ is real and reflection-even, the conjugate-balanced channels have the spatial form
\[
w_0(i)=\widehat j_0m_i,
\]
and for $r\ge1$,
\[
w_r(i)=2\widehat j_r m_i\cos\!\left(\frac{2\pi r i}{M}\right).
\]
An odd ring has $k+1$ reflection orbits. At least one is inactive, so $p\le k$. The non-vanishing lemma therefore applies to the first $p$ channel scales.

Choose one representative $i_s$ from each active reflection orbit and set
\[
x_s=\cos\!\left(\frac{2\pi i_s}{M}\right).
\]
The $x_s$ are distinct. After removal of nonzero row and column scales, the resulting $p\times p$ matrix is
\[
C_{rs}=T_r(x_s),\qquad 0\le r,s\le p-1,
\]
where $T_r$ is the Chebyshev polynomial of degree $r$. Its determinant is a nonzero constant multiple of
\[
\prod_{0\le i<j<p}(x_j-x_i),
\]
so the first $p$ channels are linearly independent. Reflection symmetry bounds the active masked space by dimension $p$, hence
\[
\boxed{\operatorname{rank}(W)=p.}
\]

\section{The stencil preserves rank}
The multiplier
\[
s_q=\cos\!\left(\frac{2\pi q}{M}\right)-1
\]
vanishes only at $q=0$, so the spatial kernel of $S$ consists exactly of constant vectors. Every vector in $\operatorname{span}(W)$ vanishes on every inactive site because each channel contains the factor $m_i$. Therefore the only constant vector in $\operatorname{span}(W)$ is zero:
\[
\ker S\cap\operatorname{span}(W)=\{0\}.
\]
Thus $S$ is injective on the channel span and
\[
\boxed{\operatorname{rank}(V)=\operatorname{rank}(SW)=\operatorname{rank}(W)=p.}
\]
This proves the theorem.

\section{Locked benchmark and computational audit}
For the locked $M=99$, $A=0.50$, $\theta=0.08$ benchmark, there are 41 active reflection orbits; the theorem gives $\operatorname{rank}(V)=41$. Under the same $A,\theta$ construction, the tested neighboring rings give 42 active reflection orbits for $M=101$ and 52 for $M=127$, hence ranks 42 and 52 respectively.

The companion verifier checks the closed shedding field, compares the simplified Fourier formulas with direct FFT coefficients for all odd $M=9,\dots,2001$, and performs exact integer/cyclotomic non-vanishing checks for $M=99,101,127$. These calculations corroborate the analytic proof but do not replace it.

\section{Claim firewall}
This theorem concerns the rank of a single-step phenomenological cyclic redistribution operator. It does not establish a microscopic physical law, a universal compression dimension, a universal value of $K_\tau$, a multi-step dynamical theorem, or a result about the separate entropy-to-horizon bridge question.

\end{document}
